
\chapter{Introduction}
\label{ch:introduction}

\section{Motivation}
\label{sec:ch1_motivation}

Small language models make it possible to run language-based applications on a
smartphone, close to the user and without a remote inference service~\citep{ondevice_lm_review}.
Running on the device is only part of an assistant's usefulness, however. An assistant that
supports repeated interactions must also retain what it learned earlier, such as a user's
preferences, earlier requests, observations of the device and the outcomes of previous
actions, and recover it when it becomes relevant again. Acting on a previously stated
preference is therefore a memory problem as well as an inference problem.

The resources for this process are limited. A smartphone shares its memory and processing
capacity with the operating system and other applications, and battery capacity and heat
limit how much computation can be sustained; studies of small-model deployment therefore
evaluate prompt processing, generation speed and memory alongside model
capability~\citep{slm_edge_deployment}. Keeping the whole interaction history in the model's
context window is the simplest way to make past information available, but the history keeps
growing, and longer inputs require more prompt processing and more storage for attention
states, depending on the architecture and runtime~\citep{ondevice_lm_review}. Even when a long
context fits in memory, processing it may take too much time or energy for repeated use on a
phone.

External memory separates the amount of information retained from the amount presented to the
model for a particular request. A system can store a long history and retrieve only a small part
of it, following the retrieval-augmented generation approach of supplying external evidence at
inference time~\citep{rag_lewis}. This creates a design choice: how should the history be
represented, which parts should be retrieved, and what should the model receive within its
available context budget? This thesis studies that choice for small language models, with
particular attention to the physical cost of memory use on a smartphone.

\section{Problem statement}
\label{sec:ch1_problem}

A useful memory strategy must make relevant past information available at a cost the
application can afford. Storage alone does not meet this requirement: the stored information
must be retrievable, the retrieved representation must preserve enough detail for the model to
use it correctly, and construction, retrieval and answer generation all add to the cost.

This thesis compares three approaches. The full-context baseline presents the dialogue
history within a configured input limit, truncating it when necessary. Vector retrieval
stores dialogue text and selects relevant turns using embedding similarity. The graph
approach asks a small language model to extract subject--relation--object triples, stores
the resulting relations and retrieves graph evidence for answering. These approaches
differ both in the work required to prepare memory and in the evidence subsequently
shown to the answer model, referred to here as the \emph{reader}.

Knowledge graphs are of particular interest because they represent explicit relations
between entities. A compact relational representation may allow the reader to receive
several relevant facts without processing their full conversational context. Its utility
depends, however, on the extraction and retrieval stages. A missing relation cannot be
recovered through traversal alone, and a short triple may omit a qualification or
temporal detail needed to answer a question. Prior comparisons of retrieval and graph
retrieval show that performance depends on the graph design and the evidence supplied
to the reader~\citep{rag_vs_graphrag}. This motivates testing specific graph-memory
choices rather than assuming that structure is beneficial by itself.

The same distinction applies to efficiency. A short query prompt does not reveal the cost
already paid to build its memory: generating a graph requires repeated model calls, while
vector indexing embeds the stored text. Whether that work pays off depends on answer quality
and on how often the memory is reused, so an on-device comparison must account for measured
physical cost as well as prompt length.

Related work already examines parts of this problem. MobileRAG evaluates retrieval for
document question answering on a smartphone, including its energy
cost~\citep{mobilerag_energy}. SmartRAG studies an on-device graph-based system and
reports memory and power measurements, while its comparison baselines use a cloud-side
simulation~\citep{smartrag}. A benchmark of memory architectures for compact mobile agents
measures latency on a laptop and lists real-phone latency and device-level cost
measurement as future work~\citep{mobile_memory_architectures}. These studies each cover part
of the problem, but they do not compare alternative memory representations for a conversation
history on one phone while measuring both answer quality and energy. This thesis makes that
comparison with small readers,
combining a broader workstation evaluation with direct smartphone measurements and a
detailed analysis of the tested graph pipeline.

\section{Research questions}
\label{sec:ch1_research_questions}

The thesis addresses three research questions.

\begin{description}
  \item[RQ1.] How do different long-term memory strategies affect the question-answering
  performance of small language models?

  \item[RQ2.] What trade-offs between answer quality and resource consumption arise when
  deploying long-term memory strategies on smartphones?

  \item[RQ3.] What are the capabilities and limitations of knowledge-graph-based memory
  for small language models under constrained context and resource budgets?
\end{description}

RQ1 establishes the quality comparison, including how it depends on retrieval depth, reader
size and the context budget. RQ2 relates answer quality measured on the phone to energy, latency, memory use
and thermal behavior; correct answers per kilojoule give a joint measure, while the individual
resource measurements explain practical deployment constraints. RQ3 asks how extraction,
evidence selection and representation affect what reaches the reader, what repairing retrieval
or restoring source text recovers, and whether Selective Provenance Expansion (SPE), a policy
developed in this thesis, can choose whether a retrieved fact is omitted, kept as a triple or
accompanied by its source text under a context budget.

The phone comparison of answer quality covers vector retrieval and truncated full context;
graph memory enters the phone study through a directly measured construction pilot. The graph
repairs and SPE are evaluated for quality on the workstation, and their prompt lengths and
retrieval times do not substitute for phone-energy measurements.

\section{Approach and scope}
\label{sec:ch1_scope}

The experiments use question answering to evaluate whether a memory system can recover
and supply information from an extended interaction history. LoCoMo provides long
conversations and questions that require using information from those
conversations~\citep{locomo}. LongMemEval evaluates long-term interactive memory across
several abilities, including information extraction, reasoning across sessions,
temporal reasoning, knowledge updates and abstention~\citep{longmemeval}. The thesis
uses specified subsets and configurations of these benchmarks, described in
Chapter~\ref{ch:methodology}.

These tasks give the comparison an observable outcome: whether the answer produced from
the available memory is correct under a common judging procedure. They also allow the
memory representation to change while the question and underlying history remain fixed.
Mobile agents motivate the study broadly, but the measured task is memory-dependent
question answering. Improvements in action selection, tool-call correctness or complete
mobile workflows would require additional evaluation.

One experimental pipeline runs in two settings that complement each other. The main
platform is a Samsung Galaxy S23 Ultra, where the memory modules run locally and each run
records the phone's answers together with their energy, battery charge, latency, memory use
and temperature. A workstation runs the same modules on more questions, five readers, both
benchmarks and several graph variants, more than a phone can process in practical time. One
automatic judge, run on the workstation outside the measured phone workload, scores both.
Device runtimes and model packaging can change the generated answers, so a workstation score
is never used as a phone score: every phone efficiency figure uses the accuracy of the answers
the phone itself produced.

The findings concern the evaluated models, histories and execution conditions. The
phone measurements come from one handset, with a smaller question sample and uneven
replication across phone conditions. Graph construction on the phone is
represented by an early pilot, not a complete graph-QA comparison. The study
also characterizes the repeatability of its automatic judge and checks its verdicts against a
blind human audit with a single annotator. These limits define how far the results can be
generalized and are developed in the relevant chapters.

\section{Contributions}
\label{sec:ch1_contributions}

The thesis contributes an empirical comparison of long-term memory strategies for small
language models and an analysis of the costs and failure mechanisms that shape their
use on a phone. Its main contributions are:

\begin{enumerate}
  \item A controlled comparison of full context, truncated and with a long window, vector retrieval
  and SLM-extracted graph memory for conversational question answering. The evaluation examines
  reader size and retrieval depth alongside memory strategy, with a second benchmark used to
  assess whether the observed ordering persists beyond the primary question set.

  \item A smartphone study that relates answer quality to measured energy and other
  device resources. It reports correct answers per unit of energy alongside latency,
  memory and thermal observations, and compares the energy of building vector and graph
  memory on the tested handset.

  \item An analysis of graph-memory limitations through successive retrieval repairs,
  source-text restoration and controls that distinguish extractor choice from reader
  choice. The analysis identifies several stages at which useful evidence can be
  lost or fail to reach the reader, rather than attributing all errors to a single
  property of knowledge graphs.

  \item The design and evaluation of Selective Provenance Expansion against criteria
  recorded before its evaluation runs. The result is negative: the tested policies did
  not reliably recover the benefit of source-text restoration under the evaluated
  budgets, and the repaired selector chose no source turns. The measured costs are
  consistent with the per-token scoring rule and the budget favouring triples, which bounds
  this per-fact policy without ruling out query-level selective expansion.
\end{enumerate}

Each of Chapters~\ref{ch:methodology} to~\ref{ch:spe} opens by stating what it adds to the work
reviewed in Chapter~\ref{ch:background}, and Tables~\ref{tab:ch9_claims} and~\ref{tab:ch9_claims_b} list the main claims
with their sample sizes, tests and status.

On the phone, vector retrieval had a higher accuracy point estimate, lower energy per
answer and shorter query time than truncated full context in every reader pair, with
5.8--26.9$\times$ more correct answers per kilojoule where the ratio is defined; on the
workstation, with full context truncated to a 4,096-token window, it had the highest accuracy point
estimate for every reader. Long-window full context, which held the whole LoCoMo history, matched or
exceeded vector retrieval on LoCoMo and trailed it, without separating, on the much longer
LongMemEval-S histories, at about 39 and 48 times the prompt tokens. Building memory in the phone graph pilot cost about 26 times more
energy per dialogue turn than vector ingestion. Restoring source text raised Gemma~4 E2B's
graph accuracy from 21.5\% to 33.0\% and narrowed the gap to vector retrieval for some readers
without closing it for any, and the repaired SPE selector chose no source turns.

\section{Thesis organization}
\label{sec:ch1_organization}

Chapter~\ref{ch:background} introduces small language models, external memory, graph
retrieval and on-device execution, and positions the thesis against related work.
Chapter~\ref{ch:methodology}
describes the shared pipeline, benchmarks, judging procedure and resource measurements.

Chapter~\ref{ch:memory_evaluation} addresses RQ1 through the memory-strategy comparison
and its sensitivity to retrieval depth and reader choice.
Chapter~\ref{ch:on_device_memory} addresses RQ2 through the smartphone measurements.
RQ3 spans Chapter~\ref{ch:graph_mechanisms}, which analyzes the graph-repair sequence,
and Chapter~\ref{ch:spe}, which presents and evaluates SPE.

Chapter~\ref{ch:discussion} brings the findings together as answers to the three research
questions and examines their limitations. Chapter~\ref{ch:conclusion} sets out the implications
for on-device memory and the remaining research directions. Supporting statistical
analyses appear in Appendix~\ref{app:on_device_statistics}, and
Appendix~\ref{app:reproducibility} records the prompts, question sets, model artifacts and
runtime settings needed to interpret and repeat the experiments.
